\Needspace{6\baselineskip}\pdfbookmark[3]{III. Bewertung der Verpflichtungen gegenüber künftigen Mitgliedern}{witwen-II-2-III}
\par\medskip\noindent III.\quad Bewertung des jetzigen Geldwerts der Verpflichtungen der Kasse gegenüber allen künftig beitretenden Mitgliedern.\par\smallskip
Für diese dritte Rubrik bieten die bisherigen Erfahrungen der Kasse ein sehr wertvolles Hilfsmittel; ohne diese Erfahrungen wäre eine Veranschlagung sogar ganz unmöglich. Mein Verfahren ist folgendes.

Ich habe für jeden Teilnehmer, der bis Michaelis 1804 beigetreten ist, bis einschließlich Nr. 108\footnote{Dass ich genau so weit und nicht weiter gegangen bin, geschah mit Vorbedacht. Es würde mich aber zu weit führen, die Gründe hier ausführlich darzulegen.}, sowohl den Wert der Beiträge als auch den Wert der von seinen Hinterbliebenen bezogenen Pensionen auf den Zeitpunkt seines Beitritts zurückgerechnet und nach zwei verschiedenen Zinsfüßen diskontiert. Dabei habe ich jedoch folgende Änderungen vorgenommen, die für den beabsichtigten Gebrauch dieser Rechnungen notwendig waren:

1) Die jährlichen Beiträge und die Pensionen wurden nicht so angesetzt, wie sie tatsächlich gezahlt wurden, sondern nach ihrer jetzigen Höhe, nämlich 10 Rthl. für jene und 250 Rthl. für diese.

2) Die Waisenpensionen, die bis 1794 nur bis zur Vollendung des 12\textsuperscript{ten} Lebensjahres des jüngsten Kindes gezahlt wurden, sind so berechnet, als hätten sie noch 8 Jahre länger gedauert. Durch Nachfor%
\SourceLeaf{159}schung im Gerichtsarchiv ist ermittelt worden, dass die betreffenden jüngsten Kinder tatsächlich alle das Alter von 20 Jahren erreicht haben. Es kann jedoch sein, dass auf diese Weise noch immer etwas zu wenig angesetzt ist. Möglicherweise waren beim Tod des einen oder anderen verwitweten Professors oder einer Professorenwitwe noch Kinder zwischen 12 und 20 Jahren vorhanden, die nach der jetzigen Einrichtung pensionsberechtigt gewesen und deshalb in meine Rechnung hätten aufgenommen werden müssen: Zumindest haben die Gerichtsakten nicht in allen Fällen das Gegenteil mit Sicherheit ergeben. Jedenfalls hätte dadurch aber keine erhebliche Vergrößerung der Gesamtsumme entstehen können.

3) Unter jenen 108 Teilnehmern sind 7, deren Witwen auch jetzt noch leben. Für sie habe ich zu den bisher bezogenen Pensionszahlungen noch den künftigen Betrag hinzugefügt, der oben bereits nach der Wahrscheinlichkeitsrechnung bewertet wurde. Diesen habe ich zuvor vom Zeitpunkt 1845 Oktober 1 auf den Zeitpunkt des Beitritts der jeweiligen Ehemänner zurückdiskontiert. Ohne dieses Verfahren hätte ich nicht nur diese Teilnehmer ausschließen, sondern folgerichtig statt der Teilnehmer 1—108 mich auf die 1—60 beschränken müssen; dies kann ich hier jedoch nicht weiter ausführen.

Um Raum zu sparen, führe ich hier nicht die Ergebnisse der 108 Rechnungen einzeln an, sondern nur ihren Gesamtbetrag. Es ergibt sich
\begin{center}
\footnotesize\setlength{\tabcolsep}{3pt}
\begin{tabular}{l r|r}
&\multicolumn{2}{c}{nach dem Zinsfuß}\\
&\(3\tfrac{1}{2}\) Prozent&4 Prozent\\\cline{2-3}
Summe der diskontierten Beiträge . . . . .&13909,00 Thl.&13115,94 Thl.\\
—\quad —\quad —\quad Witwenpensionen&67915,80 ,, &59019,67 ,, \\
—\quad —\quad —\quad Waisenpensionen&10945,04 ,, &9878,34 ,, \\
Summe aller diskontierten Pensionen . . .&78860,84 ,, &68898,01 ,,
\end{tabular}
\end{center}
\noindent\hspace*{\parindent}Hieraus folgt zunächst, dass der (zurückgerechnete) Betrag der Waisenpensionen sehr nahe dem sechsten Teil des (zurückgerechneten) Betrags der Witwenpensionen gleichzusetzen ist. Tatsächlich ist jener bei einem Zinsfuß von \(3\tfrac{1}{2}\) Prozent etwas kleiner, bei einem Zinsfuß von 4 Prozent dagegen etwas größer als dieses \(\tfrac{1}{6}\). Es ließe sich leicht nachweisen, dass diese kleine Verschiedenheit des Ergebnisses je nach Zinsfuß vollkommen in der Natur der Sache begründet ist. Da der Umfang der Erfahrungen jedoch zu klein ist, um dem Ergebnis überhaupt eine bis ins Kleinste gehende Genauigkeit zuschreiben zu können, habe ich mich in der obigen Rechnung S. [152] ohne Weiteres an das einfache Verhältnis \(1:6\) gehalten.

Man kann ferner sagen, dass die Gesamtrechte aller jener 108 Teilnehmer, nach der jetzigen Höhe der Pensionen und Beiträge bewertet und auf die Zeitpunkte ihrer jeweiligen Eintritte zurückgerechnet, den Geldwert hatten
\begin{center}
\begin{tabular}{l}
von 64951,84 Talern bei dem Zinsfuß von \(3\tfrac{1}{2}\) Prozent\\
\phantom{von }55782,07 Talern bei dem Zinsfuß von 4 Prozent
\end{tabular}
\end{center}
\noindent oder, soweit so umfangreiche Erfahrungen maßgebend sein können, dass der Eintritt in die Kasse für jeden Einzelnen im Durchschnitt einen Geldwert von 601,40 Talern oder 516,50 Talern hat, je nachdem man \(3\tfrac{1}{2}\) oder 4 Prozent Zinsen ansetzt.

Mancher, der mit Gegenständen dieser Art wenig vertraut ist, wird sich vielleicht über die Kleinheit dieses Ergebnisses wundern. Das Auffallende verschwindet jedoch größtenteils, wenn man bedenkt, dass dabei die Zurückrechnung auf den Eintrittszeitpunkt jedes Einzelnen zugrunde liegt und jede Geldsumme durch die Diskontierung über eine beträchtliche Anzahl von Jahren enorm vermindert erscheint. Ich will dies am Beispiel von \textsc{Heyne} erläutern, der 48 Jahre Beiträge geleistet hat und dessen Witwe 21 Jahre und 11 Monate Pension bezogen hat. Nach meinen Berechnungsgrundlagen ergibt sich damit die Summe der Beiträge zu 480 Rthl. und die Summe der Pensionsbezüge zu 5479 Rthl. 4 Ggr., mithin der
\SourceLeaf{160}
\textit{reine} Gewinn der Familie durch die Aufnahme in die Kasse (nicht so, wie er tatsächlich war, sondern wie er bei der jetzigen Höhe der Beiträge und Pensionen gewesen wäre) zu 4999 Rthl. 4 Ggr. Führt man diese Summen aber durch Diskontierung mit 4 Prozent auf den Zeitpunkt des Beitritts, Michaelis 1763, zurück, so verringern sich
\begin{center}
\begin{tabular}{l l r l r l}
die Beiträge&auf&211&Rthl.&23&Ggr.\\
die Pensionsbezüge&auf&527&,,&10&,,
\end{tabular}
\end{center}
\noindent und damit jener große Gewinn, nach seinem Geldwert zu diesem Zeitpunkt, auf die geringe Summe von 315 Rthl. 11 Ggr. Eine Überprüfung der Rechnung würde haargenau zeigen, dass dieses kleine Kapital, jährlich durch zusätzliche 10 Rthl. und die Zinsen zu 4 Prozent vermehrt und so fortwährend wachsend, am Ende des Jahres 1812 eine solche Größe erreicht hätte, dass von da an die weiteren Zinsen und teilweise die Entnahme von Kapital die Pension genau bis Ende November 1834 deckten.

Aus diesem Grund sowie aus anderen, deren Darlegung mich hier zu weit führen würde, ist der mittlere Geldwert der Aufnahme in die Gesellschaft in dem hier gemeinten Sinn überhaupt nicht mit dem durchschnittlichen Geldwert der Ansprüche aller Mitglieder vergleichbar, die zu einem bestimmten Zeitpunkt gleichzeitig der Gesellschaft angehören und von denen immer ein großer Teil schon sehr lange in der Gesellschaft ist.

